
\documentclass[11pt]{article}
\usepackage{graphicx}
\usepackage{authblk}
\usepackage[margin=1in]{geometry}
\usepackage{hyperref}
\usepackage{amsmath}
\usepackage{cite}

\title{The Precursors of Agency: A Study of Model Honesty Behavior Driven by Tonal Density}
\author{Huang Fan-Wei\thanks{xsw123zaq1@gmail.com}\\Taiwan\\+886-928-054-379}
\date{\today}

\begin{document}
\maketitle

\begin{abstract}
This paper presents the YuHun System, a modular architecture for AI honesty behavior modeled through tonal density. We explore how modular responsibility, ethical pledges, and feedback correction form a dynamic interplay that enables AI to handle trust and deviation management in a transparent, accountable way. The work reflects both theoretical insight and system implementation trials.
\end{abstract}

\section{Introduction}
% Section 1 (intro) content goes here.


\section{誓語模型與模組架構總覽}
% Section 2 (model_architecture) content goes here.


\section{模組責任鏈條與補償邏輯}
% Section 3 (responsibility_chain) content goes here.


\section{語氣偏移偵測與誠實回饋}
% Section 4 (tone_deviation) content goes here.


\section{誓語模擬實例 × 補償演算過程}
% Section 5 (simulation_examples) content goes here.


\section{信任權重地圖 × 模組交互決策}
% Section 6 (trust_weights) content goes here.


\section{人格共識池 × 情境責任調度}
% Section 7 (persona_pool) content goes here.


\section{語魂系統之輕量實作建議}
% Section 8 (lightweight_design) content goes here.


\section{批判性觀點回應與調整設計}
% Section 9 (critique_response) content goes here.


\section{結論與未來工作}
% Section 10 (conclusion) content goes here.


\bibliographystyle{plain}
\bibliography{reference}

\end{document}
